\section{Giới thiệu}

Trong một mạng cộng tác, hai người có nhiều đồng nghiệp chung có thể thuộc cùng nhóm nghiên cứu. Với hàng triệu liên kết, lập danh sách nhóm bằng cách kiểm tra từng cặp người là không khả thi. Phát hiện cộng đồng (community detection) khai thác trực tiếp các liên kết để tìm các vùng tương đối kết chặt. Phân hoạch thu được hỗ trợ khảo sát cấu trúc và có thể dùng làm đầu vào cho những bước phân tích tốn chi phí hơn. Khi một cộng đồng có nhiều liên kết ngoài nhóm, hoặc khi một người thuộc nhiều nhóm, bản thân cấu trúc có thể chỉ ra nhiều cách phân chia hợp lý.

Raghavan, Albert và Kumara~\cite{raghavan} đề xuất Label Propagation Algorithm (LPA). Mỗi đỉnh bắt đầu với một nhãn riêng, rồi chọn một nhãn phổ biến nhất từ láng giềng. Bản bất đồng bộ cho phép đỉnh sau đọc kết quả cập nhật của đỉnh trước trong cùng lượt quét. LPA cần quét tất cả các đỉnh và kiểm tra nhãn cuối mỗi lượt, ngay cả ở vùng đã ổn định. Các lần đổi nhãn trong thế hoà còn có thể làm một cụm lớn tràn qua ranh giới yếu. Chiến lược giữ nhãn (retention strategy), được thảo luận trong công trình Barber và Clark~\cite{barber}, tránh đổi nhãn khi nhãn hiện tại vẫn cực đại. Do đó điều kiện dừng chỉ cần một lượt không có thay đổi. Mỗi lượt vẫn chạm toàn bộ đồ thị, và việc đóng băng ranh giới sớm có thể để lại nhiều cụm nhỏ trong cùng một nhóm thực.

Một hướng khác dùng hàm chất lượng. Louvain~\cite{blondel} thực hiện cải thiện modularity cục bộ rồi gom cụm thành siêu đỉnh. Trên mạng có nút cầu nối, chuyển một đỉnh có thể khiến phần còn lại của một cộng đồng kém kết nối. Leiden thêm bước tinh chỉnh và khai thác lịch di chuyển theo hàng đợi~\cite{leiden}, cung cấp động lực kỹ thuật cho FLPA. Nếu chọn modularity làm mục tiêu thì giới hạn phân giải của chính hàm này vẫn ảnh hưởng đến các cộng đồng nhỏ~\cite{fortunato}. Chúng tôi dùng Leiden làm baseline vì bài báo gốc chạy Leiden với mục tiêu modularity.

Traag và Šubelj~\cite{flpa} đưa ra ba đóng góp. Về kỹ thuật, FLPA giữ cập nhật đa số và chỉ đánh giá lại những láng giềng có khả năng bị ảnh hưởng. Về lý thuyết, tác giả giải thích vì sao cạn hàng đợi đảm bảo nhãn cuối cực đại, rồi phân tích chi phí trên các họ đồ thị điển hình. Về thực nghiệm, tác giả so sánh bốn thuật toán trên đồ thị tổng hợp và mạng thực, khảo sát thời gian, chất lượng và độ ổn định. Phần TV1 kiểm tra các đóng góp này bằng hiện thực độc lập và một ablation kiểm soát riêng trật tự xử lý và quy tắc đưa lại đỉnh.
